\documentclass[aspectratio=169]{beamer}
\usetheme{Madrid}
\usecolortheme{whale}
\usepackage{booktabs}
\usepackage{array}
\usepackage{graphicx}
\usepackage{hyperref}
\setbeamertemplate{navigation symbols}{}

% ── Title info ────────────────────────────────────────────────
\title[SISR: CNN vs.\ Transformer]{Single Image Super-Resolution:\\
A Systematic Comparison of CNN-Based and Transformer-Based Methods}
\subtitle{Week 1 Progress Report}
\author{Student Name}
\institute{M.Sc.\ MISCADA · Durham University}
\date{May 2026}

\begin{document}

% ── Slide 1: Title ────────────────────────────────────────────
\begin{frame}
  \titlepage
\end{frame}

% ── Slide 2: Outline ─────────────────────────────────────────
\begin{frame}{Outline}
  \tableofcontents
\end{frame}

\section{Motivation \& Research Question}

% ── Slide 3: Motivation ───────────────────────────────────────
\begin{frame}{Motivation}
  \begin{columns}
    \column{0.55\textwidth}
      \textbf{Single Image Super-Resolution (SISR)} recovers a high-resolution (HR) image from a single low-resolution (LR) input.\\[6pt]
      \textbf{Why it matters:}
      \begin{itemize}
        \item Satellite \& aerial imaging
        \item Medical diagnostics (MRI, CT)
        \item Surveillance and forensics
        \item Consumer photography
      \end{itemize}
      \vspace{6pt}
      \textbf{The challenge:} The problem is inherently \emph{ill-posed} — infinitely many HR images can correspond to the same LR input.
    \column{0.42\textwidth}
      \centering
      \textit{Bicubic (baseline)}\\[4pt]
      $\downarrow$\\[2pt]
      \textit{Deep-learning SR}\\[4pt]
      $\downarrow$\\[2pt]
      \textit{Sharper edges, finer textures}
  \end{columns}
\end{frame}

% ── Slide 4: Research Question ────────────────────────────────
\begin{frame}{Research Question}
  \begin{block}{Core Question}
    How do \textbf{5 landmark CNN-based} SISR models (SRCNN, FSRCNN, VDSR, ESPCN, EDSR) and \textbf{5 leading Transformer-based} models (RCAN, IPT, SwinIR, HAT, DRCT) compare when trained and evaluated under \emph{identical} controlled conditions, and what architectural factors most significantly impact performance, efficiency, and generalisation?
  \end{block}
  \vspace{8pt}
  \textbf{What makes this novel:}
  \begin{itemize}
    \item Existing papers report results under \emph{different} training settings
    \item No systematic head-to-head comparison on a unified benchmark suite exists
    \item This study enforces a single training dataset, optimiser, and scale factor
  \end{itemize}
\end{frame}

\section{Literature Review}

% ── Slide 5: CNN-Based Methods ────────────────────────────────
\begin{frame}{CNN-Based SISR Methods (2014–2017)}
  \begin{center}
  \small
  \begin{tabular}{@{}lllll@{}}
    \toprule
    \textbf{Model} & \textbf{Year} & \textbf{Params} & \textbf{Key Innovation} & \textbf{Set5 PSNR $\times$4} \\
    \midrule
    Bicubic   & --   & --    & Fixed interpolation            & 28.42 dB \\
    SRCNN     & 2014 & 57K   & First end-to-end CNN for SR    & 28.95 dB \\
    FSRCNN    & 2016 & 12K   & Post-upsampling; 56$\times$ faster & 29.21 dB \\
    VDSR      & 2016 & 665K  & 20-layer global residual       & 29.77 dB \\
    ESPCN     & 2016 & 20K   & Sub-pixel convolution          & 29.55 dB \\
    EDSR      & 2017 & 43M   & Removes BatchNorm; NTIRE 2017  & 30.52 dB \\
    \bottomrule
  \end{tabular}
  \end{center}
  \vspace{4pt}
  \begin{itemize}
    \item Steady PSNR gain of \textbf{+2.1 dB} from Bicubic to EDSR over three years
    \item Key trade-off: accuracy (EDSR) vs.\ speed (FSRCNN, ESPCN)
  \end{itemize}
\end{frame}

% ── Slide 6: Transformer-Based Methods ───────────────────────
\begin{frame}{Transformer-Based SISR Methods (2018–2024)}
  \begin{center}
  \small
  \begin{tabular}{@{}lllll@{}}
    \toprule
    \textbf{Model} & \textbf{Year} & \textbf{Params} & \textbf{Key Innovation} & \textbf{Set5 PSNR $\times$4} \\
    \midrule
    RCAN   & 2018 & 15.6M & Channel attention (RIR)          & 30.77 dB \\
    IPT    & 2021 & 115M  & Large-scale pre-trained ViT      & 32.64 dB \\
    SwinIR & 2021 & 11.8M & Shifted-window self-attention    & 32.92 dB \\
    HAT    & 2023 & 20.8M & Hybrid + overlapping cross-attn  & 33.04 dB \\
    DRCT   & 2024 & 14.1M & Dense-residual connections       & 33.11 dB \\
    \bottomrule
  \end{tabular}
  \end{center}
  \vspace{4pt}
  \begin{itemize}
    \item Transformers achieve \textbf{+2.6 dB} over best CNN (EDSR $\to$ DRCT)
    \item SwinIR's shifted-window attention is the dominant backbone paradigm
    \item Efficiency improving: DRCT (14.1M) outperforms IPT (115M)
  \end{itemize}
\end{frame}

\section{Experimental Setup}

% ── Slide 7: Datasets \& Metrics ─────────────────────────────
\begin{frame}{Benchmark Datasets \& Evaluation Metrics}
  \begin{columns}
    \column{0.5\textwidth}
      \textbf{Training:}
      \begin{itemize}
        \item DIV2K — 800 HR images, diverse content
      \end{itemize}
      \vspace{6pt}
      \textbf{Test Benchmarks:}
      \begin{itemize}
        \item \textbf{Set5} (5 imgs) — quick sanity check
        \item \textbf{Set14} (14 imgs) — standard evaluation
        \item \textbf{BSD100} (100 imgs) — generalisation
        \item \textbf{Urban100} (100 imgs) — textures \& structure (hardest)
      \end{itemize}
    \column{0.47\textwidth}
      \textbf{Metrics:}
      \begin{itemize}
        \item \textbf{PSNR} (dB) — pixel-level fidelity\\
              $10\log_{10}(\text{MAX}^2/\text{MSE})$
        \item \textbf{SSIM} — structural similarity\\
              (luminance, contrast, structure)
        \item Both measured on \textbf{Y channel} (YCbCr)
        \item \textbf{Inference time} (ms/image) on a fixed GPU
      \end{itemize}
  \end{columns}
\end{frame}

\section{Research Gaps}

% ── Slide 8: Research Gaps ────────────────────────────────────
\begin{frame}{Identified Research Gaps}
  \begin{enumerate}
    \item \textbf{Efficiency–accuracy trade-off is under-studied}\\
      Large models (IPT 115M, EDSR 43M) are impractical for edge devices.

    \vspace{6pt}
    \item \textbf{Benchmark saturation}\\
      Set5 and Set14 are too small and too idealised; BSD100 and Urban100 are better but still assume bicubic degradation.

    \vspace{6pt}
    \item \textbf{Synthetic degradation gap}\\
      Nearly all models are trained on bicubic downsampling, not real-world noise, JPEG compression, or camera blur.

    \vspace{6pt}
    \item \textbf{No fair unified comparison}\\
      Published results use different training datasets, optimisers, and scale factors — direct comparison is misleading.
  \end{enumerate}
\end{frame}

\section{Project Plan}

% ── Slide 9: Methodology ─────────────────────────────────────
\begin{frame}{Experimental Methodology}
  \begin{block}{Controlled Comparison Protocol}
    All 10 models share: \textbf{same training data (DIV2K)} · \textbf{same optimiser} · \textbf{same scale factor ($\times$4)} · \textbf{same evaluation pipeline}
  \end{block}
  \vspace{8pt}
  \begin{columns}
    \column{0.48\textwidth}
      \textbf{CNN models} (train from scratch):\\
      SRCNN · FSRCNN · VDSR · ESPCN · EDSR
    \column{0.48\textwidth}
      \textbf{Transformer models} (pre-trained + fine-tune):\\
      RCAN · IPT · SwinIR · HAT · DRCT
  \end{columns}
  \vspace{10pt}
  \textbf{Planned ablations:}
  \begin{itemize}
    \item Depth / channel width impact on PSNR
    \item Attention type (channel vs.\ window vs.\ cross)
    \item Parameter count vs.\ reconstruction quality scatter plot
  \end{itemize}
\end{frame}

% ── Slide 10: Timeline & Next Steps ──────────────────────────
\begin{frame}{Project Timeline \& Next Steps}
  \begin{center}
  \small
  \begin{tabular}{@{}llp{6.5cm}@{}}
    \toprule
    \textbf{Phase} & \textbf{Period} & \textbf{Deliverable} \\
    \midrule
    1 — CNN Training     & 27 Apr – 18 May & SRCNN, FSRCNN, VDSR, ESPCN, EDSR evaluated on Set5/Set14/BSD100 \\
    2 — Transformer      & 19 May – 15 Jun & RCAN, IPT, SwinIR, HAT, DRCT; extend to Urban100 \\
    3 — Ablation         & 16 Jun – 29 Jun & Ablation studies, qualitative analysis, efficiency profiling \\
    4 — Results          & 30 Jun – 13 Jul & Full comparison tables, threat-to-validity analysis \\
    5 — Write-up         & 14 Jul – 4 Sep  & Dissertation draft, supervisor review, final submission \\
    \bottomrule
  \end{tabular}
  \end{center}
  \vspace{8pt}
  \textbf{Immediate next steps (this week):}
  \begin{itemize}
    \item Finalise data pipeline and training scripts
    \item Complete SRCNN and FSRCNN training runs
    \item Validate evaluation metrics against published numbers ($\pm$0.1 dB)
  \end{itemize}
\end{frame}

\end{document}
